\documentclass{rvcepaper}
\begin{document}

% ---------- HS351TA Entrepreneurship & IPR (CIE - II, 16 Nov 2026)
\begin{iempaper}
\paperinfo{dept={INDUSTRIAL ENGINEERING \& MANAGEMENT},date={16\textsuperscript{th} November 2026},exam={CIE -- II},maxmarks={50},sem={V},prog={UG},duration={90 Min},
 title={ENTREPRENEURSHIP \& INTELLECTUAL PROPERTY RIGHTS},code={HS351TA}}
\iemhead
\begin{cietablemc}{M}{CO}
\qmc{1.}{A Bengaluru-based startup markets packaged fruit juice under the name ``FreshDrop''. An established company already sells similar products under the mark ``FreshDrip''. Both brands use similar colours, lettering and packaging. Explain the concept, functions and different forms of trademarks. Analyse whether the two marks may be considered deceptively similar, stating the relevant factors and suitable case law. Also discuss the conditions for registration and the procedure for registering a trademark in India.}{10}{4}
\qmc{2.}{Discuss the absolute grounds for refusal of registration and the categories of non-registrable marks with suitable examples and case studies.}{10}{4}
\qmc{3.}{A company introduces an electronic product under the brand name ``ElectroMax,'' while a competitor markets a similar product under the name ``ElectroMaxx'' with a similar logo. Analyse the concept of deceptive similarity, the procedure for trademark registration, and the legal remedies available for trademark infringement.}{10}{3}
\qmc{4.}{Discuss how entrepreneurs utilize break-even analysis as a tool for risk management and financial planning.}{10}{2}
\qmc{5.}{Define SWOC Analysis and explain its four key components (Strengths, Weaknesses, Opportunities, and Challenges) in the context of strategic planning. Discuss how entrepreneurs use SWOC Analysis alongside Vision, Mission, Goals, and Objectives to formulate an effective business strategy.}{10}{2}
\end{cietablemc}
\end{iempaper}

% ---------- IM352IA Operations Management (CIE - II, 16 Nov 2026)
\begin{iempaper}
\paperinfo{dept={INDUSTRIAL ENGINEERING \& MANAGEMENT},date={16\textsuperscript{th} November 2026},exam={CIE -- II},maxmarks={50},sem={V},prog={UG},duration={90 Min},
 title={Operations Management},code={IM352IA}}
\iemhead
\begin{cietablemc}{M}{CO}
\qmc{1.}{Explain the systematic procedure for capacity planning and discuss the various stages involved in formulating an effective capacity plan.}{10}{4}
\qmc{2.}{A company is operating at 90\% of its existing production capacity and is experiencing frequent delays in meeting customer orders. Analyze the situation and suggest suitable long-term capacity planning measures.}{10}{3}
\qmc{3.}{Explain the different components of a time series forecast. Differentiate between qualitative and quantitative forecasting methods.}{10}{3}
\qmc{4.}{You have been asked to put together a capacity plan for a critical operation at the Surefoot Sandal Company. Your capacity measure is number of machines. Three products are manufactured. The time standards, lot sizes, and demand forecasts are given in the following table. The firm operates two 8-hour shifts, 5 days per week, 50 weeks per year. Experience shows that a capacity cushion of 5 percent is sufficient.
{\small\qtabc{|l|C{2.4cm}|C{2.4cm}|C{2.3cm}|C{2.7cm}|}{\hline & \multicolumn{2}{c|}{\textbf{Time standards}} & & \\\cline{2-3}
\textbf{Sandals} & \textbf{Processing time (hr/pair)} & \textbf{Set up time (hr/pair)} & \textbf{Lot size (pairs/lot)} & \textbf{Demand forecast (pairs/yr)}\\\hline
Men's & 0.05 & 0.5 & 240 & 80000\\\hline
Women's & 0.1 & 2.2 & 180 & 60000\\\hline
Children's & 0.02 & 3.8 & 360 & 120000\\\hline}}
\begin{romi}
\item If the operation currently has two machines, what is the capacity gap?
\item How many machines are needed?
\end{romi}}{10}{4}
\qmc{5.}{The actual demand and forecast demand for six months are given below:
\qtabc{|l|c|c|c|c|c|c|}{\hline \textbf{Month} & 1 & 2 & 3 & 4 & 5 & 6\\\hline
Actual Demand & 500 & 520 & 550 & 580 & 600 & 620\\\hline
Forecast Demand & 480 & 530 & 540 & 560 & 590 & 610\\\hline}
Calculate:
\begin{romi}
\item Cumulative Forecast Error (CFE)
\item Mean Forecast Error (MFE)
\item Mean Absolute Deviation (MAD)
\item Tracking Signal
\end{romi}}{10}{3}
\end{cietablemc}
\end{iempaper}

\end{document}
